% Seção 2 — Contexto institucional. Condensada de paper/secoes/02-background.tex
% e docs/legislacao/README.md. Fatos de 2011: docs/ars/19 §2.3 (a conferir).

\chapter{Contexto institucional}
\label{cap:contexto}

\section{A Lei Zé Maria do Tomé}

A Lei Estadual nº~16.820, sancionada em 8 de janeiro de 2019 e vigente desde o
dia seguinte \cite{ceara2019lei16820}, acrescentou o art.~28-B à Lei Estadual
nº~12.228/1993 \cite{ceara1993lei12228}, a norma cearense sobre agrotóxicos. O caput veda ``a pulverização aérea de agrotóxicos
na agricultura no Estado do Ceará''; o §~2º estende a vedação à dispersão por
aeronave para controle vetorial. Dois traços da redação importam ao desenho. A
proibição é de \textbf{método}: nenhum princípio ativo é nomeado, e o produto
continua disponível para aplicação terrestre. E a proibição é \textbf{estadual
e simultânea}: alcança os 184 municípios na mesma data.

O Quadro~\ref{qua:cronologia} resume a cronologia. Três marcos antecedem o
tratamento e são tratados separadamente: a \emph{notícia} (apresentação do projeto de lei que deu origem à norma, em
2015, precedida de proposição semelhante em 2011), a \emph{certeza} (aprovação unânime, em dezembro de
2018) e a \emph{obrigação} (vigência, em janeiro de 2019). Em dezembro de 2024,
a Lei nº~19.135 \cite{ceara2024lei19135} deu nova redação ao artigo e abriu exceção para aeronaves
remotamente pilotadas. O estado não revogou a proibição, mas trocou o
instrumento: de cota zero para regulação de desempenho. Por isso a janela de
análise termina em 19 de dezembro de 2024.

\begin{quadro}[htb]
\caption{\label{qua:cronologia}Cronologia legal do tratamento}
\centering
\small
\begin{tabular}{@{}lp{11.2cm}@{}}
\toprule
Data & Evento \\
\midrule
09/12/1993 & Lei Estadual nº 12.228: norma-mãe sobre agrotóxicos no Ceará. \\
20/11/2009 & Limoeiro do Norte proíbe a pulverização aérea por lei municipal (Lei nº 1.478). \\
2010 & Assassinato de Zé Maria do Tomé (abril). Em maio, a lei municipal é revogada (Lei nº 1.511). \\
2011 & Notícia de proibição da prática na Chapada do Apodi (março); ação civil pública do Ministério Público e projeto de lei estadual com o mesmo objeto (julho). Alcance, desfecho e tramitação a conferir. \\
24/02/2015 & Apresentação do Projeto de Lei nº 18/2015 (marco de notícia). \\
18/12/2018 & Aprovação por unanimidade na Assembleia Legislativa (marco de certeza). \\
09/01/2019 & Vigência da Lei nº 16.820 (marco de obrigação). \\
05/2023 & O Supremo Tribunal Federal julga a lei constitucional (ADI 6137). \\
19/12/2024 & Lei nº 19.135: exceção para drones. Fim da janela de análise. \\
\bottomrule
\end{tabular}
\fonte{Elaboração própria, a partir dos textos legais e da tramitação legislativa; os eventos de 2009 a 2011 provêm de fontes secundárias, a conferir.}
\end{quadro}

\section{A Chapada do Apodi, a cultura e o químico}

A Chapada do Apodi, no Baixo Jaguaribe, concentra a fruticultura irrigada do
Ceará, com destaque para banana e melão. É a região à qual a lei responde e
onde a exposição da população à deriva foi documentada antes da proibição. Na
bananicultura, a aplicação aérea serve ao controle de doenças fúngicas da
folha, com aplicações repetidas ao longo do ano. No análogo equatoriano, elas
se intensificam na estação chuvosa, quando a pressão da doença aumenta
\cite{calzada2023bananas}; se o mesmo vale para o Ceará é questão que a etapa
de medição deste projeto verifica.

Esse perfil distingue o caso da literatura de referência, construída em grande
parte sobre glifosato em grãos transgênicos. A tabela de análises de água
reproduzida na justificativa do Projeto de Lei nº~18/2015, atribuída ao Dossiê
ABRASCO \cite{ceara2015pl18}, registra 23 pontos de coleta na Chapada em 2009: o fungicida
procimidona e o inseticida carbaril aparecem em todos, carbofurano em 18,
fenitrotiona em 16 e glifosato em apenas 4. Trata-se de fungicida e inseticida
sobre fruta irrigada. Os desenhos de referência transferem na forma, não no
mecanismo, e o análogo empírico adequado é o da fumigação aérea de bananais
\cite{calzada2023bananas}.

\section{O que se sabe sobre onde havia avião}
\label{sec:onde-aviao}

O dado público mais direto sobre o uso de aeronaves é o Censo Agropecuário de
2006, que registra os estabelecimentos que aplicaram agrotóxicos por avião. No Ceará eram 36 estabelecimentos, em 7 dos 184 municípios
\cite{ibge2006censoagro}. A prática era, portanto, geograficamente concentrada,
ao passo que a bananicultura está presente na maior parte do estado: a maioria
dos municípios com a cultura não tinha registro de avião.

Dois fatos posteriores ao Censo impedem tomar esse retrato como definitivo. Em
março de 2011, noticiou-se que as empresas de fruticultura da Chapada estariam
proibidas de pulverizar por aeronave, e em julho do mesmo ano o Ministério
Público ajuizou ação civil pública com esse pedido. O alcance da primeira
medida e o desfecho da segunda não foram ainda conferidos em fonte primária.
Saber se a prática persistiu até 2018 é condição para saber o que a lei de
2019 removeu, e por isso a etapa de medição do
Capítulo~\ref{cap:metodologia} antecede a estimação principal.
